\chapter{编译器指令}
\label{ch:23}

第 19 章的 \texttt{\%} 改的是源码文本：一段代码要不要参与编译、几个片段怎么拼。第 20 章的 \texttt{@} 贴的是运行期资料：程序自己能把标签读回来。这一章讲第三种编译期手段：编译器指令——以 \texttt{\#} 开头、写在声明前面的那几行。

一句话概括：编译器指令 = 给声明贴的、只有编译器看的一张便条。

它不改文本，也不留到运行时。它的作用只有一个：让编译器在处理这个声明时多知道一条规矩——这个类型不许复制、这些样板方法由你生成、这个运算符别再自动推导。

\section{三种编译期手段的分工}

三种都以特殊符号开头，很容易记混。差别其实很清楚：谁在读它、什么时候读。

{\small\setlength{\tabcolsep}{4pt}
\begin{longtable}{>{\raggedright\arraybackslash}p{\dimexpr 0.462\linewidth-2\tabcolsep\relax}>{\raggedright\arraybackslash}p{\dimexpr 0.080\linewidth-2\tabcolsep\relax}>{\raggedright\arraybackslash}p{\dimexpr 0.104\linewidth-2\tabcolsep\relax}>{\raggedright\arraybackslash}p{\dimexpr 0.143\linewidth-2\tabcolsep\relax}>{\raggedright\arraybackslash}p{\dimexpr 0.191\linewidth-2\tabcolsep\relax}}
\toprule
写法 & 名字 & 谁读 & 什么时候 & 典型用途 \\
\midrule
\endhead
\texttt{\%define} / \texttt{\%if} & 预处理 & 预处理器 & 编译前，只认文本 & 条件编译、开关日志 \\
\texttt{\#derive} / \texttt{\#Non\allowbreak{}Copyable} / \texttt{\#No…Replace} / \texttt{\#No\allowbreak{}Right…} & 编译器指令 & 词法 + 语义层 & 编译期，认声明结构 & 改声明语义、关闭自动推导 \\
\texttt{@Route(\allowbreak{}.\allowbreak{}.\allowbreak{}.\allowbreak{})\allowbreak{}} & 注解 & 运行期反射 & 运行时 & 框架读的元数据 \\
\bottomrule
\end{longtable}
}

三者的代价也不同：

\begin{itemize}

\item 预处理最"脏"——它看不见类型，只看见字；

\item 编译器指令最"准"——它挂在具体声明上，接口名写错、位置写错都会在编译期报错；

\item 注解最"轻"——不影响语义，但需要反射能力才能读到。

\end{itemize}

本章只讲中间那一行。ChaLi 目前有四个编译器指令：

\begin{center}
\begin{tabular}{ll}
\toprule
指令 & 作用 \\
\midrule
\texttt{\#Non\allowbreak{}Copyable} & 这个 \texttt{opcode} 类型不许复制 \\
\texttt{\#derive(\allowbreak{}Equal,\allowbreak{} Hash,\allowbreak{} Clone)\allowbreak{}} & 按字段生成接口方法 \\
\texttt{\#No<\allowbreak{}接口>\allowbreak{}Replace} & 这个类型的该运算符不再自动推导 \\
\texttt{\#No\allowbreak{}Right<\allowbreak{}接口>\allowbreak{}} & 只关掉"右操作数当接收者"那一条推导 \\
\bottomrule
\end{tabular}
\end{center}

\section{\texttt{\#Non\allowbreak{}Copyable}：这个类型不许复制}

有的类型天生不该有第二份。文件句柄、进程、socket、数据库连接——拷贝一个句柄，就是两份所有权、两次 \texttt{close(\allowbreak{})\allowbreak{}}，系统里立刻多出一个悬空资源。

\texttt{opcode} 声明把这种宿主对象交给 C 侧管理，其中"能不能复制"由声明体里有没有绑定\texttt{copy} 决定。不绑定 \texttt{copy}，就写成 \texttt{\#Non\allowbreak{}Copyable}：

\par
\medskip
\begin{lstlisting}[style=chali]
opcode Number {
    new = createNumber;
    drop = destroyNumber;
    copy = copyNumber;
};

#NonCopyable
opcode FileHandle {
    new = openFile;
    drop = closeFile;
};
\end{lstlisting}
\par
\medskip

\texttt{Number} 可以随便复制；\texttt{File\allowbreak{}Handle} 不行——复制初始化、赋值、按值传参、返回时复制全部报错。想转移所有权，用 \texttt{move(\allowbreak{}\&x)\allowbreak{}}：

\par
\medskip
\begin{lstlisting}[style=chali]
FileHandle handle = FileHandle::new("log.txt");
FileHandle moved = move(&handle);      // 可以：转移
// FileHandle copied = handle;         // 报错：不可复制的值禁止复制
\end{lstlisting}
\par
\medskip

诊断会说清两种来路："实现 Closeable 或声明了 \#NonCopyable"——因为实现 \texttt{Closeable} 的结构体自动获得这条规则，不需要再写指令。第 8 章说过，\texttt{Closeable} 意味着"出块时自动关闭"；一个值被复制成两份，\texttt{close(\allowbreak{})\allowbreak{}} 就会跑两遍，所以语言直接把复制堵死：

\par
\medskip
\begin{lstlisting}[style=chali]
struct Handle : Closeable
{
    int64 id;

    void close()
    {
        println("close {}", self.id);
    }
};

void consume(Handle handle)
{
    println("consume {}", handle.id);
}

void main()
{
    Handle first;
    first.id = 1;
    consume(move(&first));

    Handle second;
    second.id = 2;
    Handle third = move(&second);
    println("third={}", third.id);


}
\end{lstlisting}
\par
\medskip

\par
\medskip
\begin{lstlisting}[style=out]
consume 1
close 1
third=2
close 2
\end{lstlisting}
\par
\medskip

注意输出里没有 \texttt{close 1} 之后的自动关闭：\texttt{first} 已经把所有权交出去了，出块时它不再是拥有者。\texttt{third} 是移动过来的新绑定，它才负责收尾。

要不要写 \texttt{move} 的判断很朴素：谁最后拿着这个值，谁负责关。转移写 \texttt{move(\allowbreak{}\&x)\allowbreak{}}，只是借用就传 \texttt{T mut\&} / \texttt{T const\&}，函数内部造出来的局部值直接 \texttt{return}。

\section{\texttt{\#derive(\allowbreak{}Equal,\allowbreak{} Hash,\allowbreak{} Clone)\allowbreak{}}：字段级样板交给编译器}

写值语义的类型时，\texttt{=\allowbreak{}=\allowbreak{}}、\texttt{hash(\allowbreak{})\allowbreak{}}、\texttt{clone(\allowbreak{})\allowbreak{}} 三件套几乎是机械劳动：逐字段比较、逐字段混合哈希、逐字段复制。字段一多，漏一个就出 bug，而且很难测出来。

\texttt{\#derive} 让编译器按字段生成：

\par
\medskip
\begin{lstlisting}[style=chali]
#derive(Equal, Hash, Clone)
struct Point
{
    int64 x;
    int64 y;
    String label;
};

void main()
{
    Point origin;
    origin.x = 0;
    origin.y = 0;
    origin.label = "O";

    Point copy = origin.clone();
    println("equal={}", origin == copy);
    println("sameHash={}", origin.hash(origin) == copy.hash(copy));

    copy.label = "O2";
    println("afterChange={}", origin == copy);
    println("original={}", origin.label);

}
\end{lstlisting}
\par
\medskip

\par
\medskip
\begin{lstlisting}[style=out]
equal=1
sameHash=1
afterChange=0
original=O
\end{lstlisting}
\par
\medskip

（布尔值用 \texttt{\{\}} 打印出来是 \texttt{1} / \texttt{0}，与 C 的习惯一致。）

\texttt{clone(\allowbreak{})\allowbreak{}} 是深拷贝：改副本的 \texttt{label} 不影响原件。\texttt{hash(\allowbreak{})\allowbreak{}} 保证"相等的值哈希相同"，所以派生出来的类型可以直接当 \texttt{Map} / \texttt{Set} 的键。

三条硬规矩：

\begin{enumerate}

\item 白名单只有 \texttt{Equal}、\texttt{Hash}、\texttt{Clone}。 写别的名字直接报错。

\item 不能派生 \texttt{Drop} 或 \texttt{Closeable}。 这两个是析构协议，值语义已经按字段自动递归清理了，再派生一次就变成"析构两遍"。

\item 嵌套字段必须自己也有协议。 字段是自定义结构体时，那个结构体得实现对应接口，或者自己也 \texttt{\#derive}。引用、函数值、裸数组这些没有相等 / 哈希语义的字段会被拒绝——编译器不会按内存字节偷偷比较。

\end{enumerate}

派生方法的生成分两步：语义层从字段元数据生成受限的方法源码，交给现有的宏模板与生成解析器，产出的普通方法再走一遍完整的类型、借用与降级检查。也就是说，派生不是特权通道，它生成的方法和手写的一样受检查。

\section{\texttt{\#No\allowbreak{}Add\allowbreak{}Replace}：关掉运算符的自动推导}

第 9 章提过运算符重载：结构体继承 \texttt{Add} 这类左接口，右操作数那侧要自己负责就继承\texttt{Right\allowbreak{}Add}。为了少写代码，编译器还替你做三类推导：

\begin{itemize}

\item 交换律换边：\texttt{+ * =\allowbreak{}=\allowbreak{} !=\allowbreak{} \& | \textasciicircum{}} 只写了左接口时，\texttt{5 + a} 自动换边借 \texttt{a + 5}；

\item 比较对偶：\texttt{5 >\allowbreak{} a} 等价于 \texttt{a <\allowbreak{} 5}；

\item 取反：\texttt{a >\allowbreak{}=\allowbreak{} 5} 就是 \texttt{!(\allowbreak{}a <\allowbreak{} 5)\allowbreak{}}。

\end{itemize}

绝大多数时候这很省事。但有些类型的"自动换边"是错的：比如一个表示"链式追加"的类型，\texttt{a + 5} 是"把 5 接到 a 后面"，而 \texttt{5 + a} 大概该是"把 a 接到 5 后面"——两种语义根本不同，推导出来的那一份只是"看起来能编"。

这时在 \texttt{struct} 前写 \texttt{\#No\allowbreak{}Add\allowbreak{}Replace}：

\par
\medskip
\begin{lstlisting}[style=chali]
#NoAddReplace
struct Vec : Add<int64, int64>
{
    int64 v;

    int64 leftAdd(int64 other)
    {
        return self.v + other;
    }
};

void main()
{
    Vec a;
    a.v = 10;
    println("left={}", a + 5);

}
\end{lstlisting}
\par
\medskip

\par
\medskip
\begin{lstlisting}[style=out]
left=15
\end{lstlisting}
\par
\medskip

\texttt{\#No\allowbreak{}Add\allowbreak{}Replace} 里的接口名就是左接口名（\texttt{Add} / \texttt{Multiply} / \texttt{Less} / \texttt{Less\allowbreak{}Or\allowbreak{}Equal} …），指令本身不带括号。写完之后，\texttt{a + 5} 照常调用 \texttt{left\allowbreak{}Add}，但 \texttt{5 + a} 不再换边：

\par
\medskip
\begin{lstlisting}[style=out]
chali: samples-error/operator_no_add_replace.chl 第 16 行第 5 列：运算符 + 的自动推导被 #NoAddReplace 关闭：int64 + Vec（直接写出与实现匹配的运算符，或显式实现 RightAdd<R,T> 这类右接口）
\end{lstlisting}
\par
\medskip

报错文案会点出是哪一条指令挡的，并给出两条出路：自己实现 \texttt{Right\allowbreak{}Add<\allowbreak{}int64,\allowbreak{}int64>\allowbreak{}}，或者接受"这个方向不支持"。

指令只关推导，不关直接写法：左向 \texttt{a + 5} 与显式实现的 \texttt{Right\allowbreak{}Add} 都不受影响。这条边界很重要——它保证"加一条指令"不会把本来能编的代码弄坏，只会把"靠猜"的那部分收回去。

\section{\texttt{\#No\allowbreak{}Right\allowbreak{}Add}：只关右向}

\texttt{\#No\allowbreak{}Add\allowbreak{}Replace} 是"这个接口的所有推导都别做了"。更多时候你只想关一半：左向的 \texttt{a + 5} 是对的，右向的 \texttt{5 + a} 是错的。这时候用 \texttt{\#No\allowbreak{}Right<\allowbreak{}接口>\allowbreak{}}：

\par
\medskip
\begin{lstlisting}[style=chali]
#NoRightAdd
struct Chain : Add<int64, int64>
{
    int64 v;

    int64 leftAdd(int64 other)
    {
        return self.v + other;
    }
};

void main()
{
    Chain chain;
    chain.v = 7;
    println("chain + 1 = {}", chain + 1);

}
\end{lstlisting}
\par
\medskip

\par
\medskip
\begin{lstlisting}[style=out]
chain + 1 = 8
\end{lstlisting}
\par
\medskip

而 \texttt{1 + chain} 会报错：

\par
\medskip
\begin{lstlisting}[style=out]
chali: samples-error/operator_right_forbidden.chl 第 13 行第 5 列：右向被编译器指令关闭：Vec 声明了 #NoRightAdd，int64 + Vec 不再推导；要去掉该指令，或改为实现 RightAdd<R,T>
\end{lstlisting}
\par
\medskip

两个指令的分工，用一张表说清：

\begin{center}
\begin{tabular}{llll}
\toprule
写法 & 左向 \texttt{a + 5} & 右向 \texttt{5 + a} & 适用场合 \\
\midrule
什么都不写 & 直接实现 & 交换律换边 & 两个方向确实等价 \\
\texttt{Right\allowbreak{}Add} & 直接实现 & 走 \texttt{right\allowbreak{}Add} & 右向语义不同，自己实现 \\
\texttt{\#No\allowbreak{}Right\allowbreak{}Add} & 直接实现 & 报错 & 右向没有意义，明确拒绝 \\
\texttt{\#No\allowbreak{}Add\allowbreak{}Replace} & 直接实现 & 报错，对偶 / 取反也不推 & 这个接口的推导整体不可信 \\
\bottomrule
\end{tabular}
\end{center}

值得强调一段历史：\texttt{\#No\allowbreak{}Right\allowbreak{}Add} 这类开关早期是空标记接口，写成\texttt{struct Chain :\allowbreak{} Add<\allowbreak{}int64,\allowbreak{}int64>\allowbreak{},\allowbreak{} No\allowbreak{}Right\allowbreak{}Add}。之所以搬成指令，是因为"禁止某条推导" 根本不涉及方法契约——它不需要被实例化、不需要类型实参、也不会被调用，把它混进继承列表会让人误以为它是个能力。声明级的东西就该写在声明前面。

\section{指令会在哪里被检查}

编译器指令看起来只是"贴张条"，但每一条都在语义层被认真核对：

\begin{center}
\begin{tabular}{ll}
\toprule
检查 & 说明 \\
\midrule
只能修饰声明 & \texttt{\#derive} 只能贴 struct / enum；运算符开关只能贴 struct \\
接口名必须存在 & \texttt{\#No\allowbreak{}Add\allowbreak{}Replac} 这种笔误直接报错，不会静默忽略 \\
接口必须有可关的东西 & \texttt{\#No\allowbreak{}Subtract\allowbreak{}Replace} 报错——\texttt{-} 本来就没有自动推导 \\
不能重复 & 同一接口写两遍报错 \\
泛型模板会带走指令 & 模板声明前的指令随实例化贴回每个具体实例 \\
\bottomrule
\end{tabular}
\end{center}

最后一条要展开说一句。写泛型结构体时，指令是写在模板上的：

\par
\medskip
\begin{lstlisting}[style=chali]
#NoAddReplace
struct Wrapper<T> { T value; };
\end{lstlisting}
\par
\medskip

单态化会在编译期为 \texttt{Wrapper<\allowbreak{}int64>\allowbreak{}}、\texttt{Wrapper<\allowbreak{}String>\allowbreak{}} 各生成一个实例。实例是从模板文本展开出来的——如果指令留在原地，实例就丢了这条规矩。所以收集模板时会把指令先摘走，贴在实例声明的前面，和 \texttt{\#derive} 一样处理。你在源码里写一次，所有实例都守这条规矩。

\section{什么时候该用指令}

指令不是"更高级的语法"，它解决的是声明层面说不清的话。判断标准可以压成两条：

\begin{itemize}

\item 这件事属于类型本身，不属于任何一次调用。 "这个类型不许复制"、"这个类型的方向加法语义不对"——是类型的事实，写在类型上；

\item 编译器不写这一行就会猜错。 派生三件套、运算符推导都属于"编译器替你做决定"，指令的作用就是把这些决定的开关交回你手里。

\end{itemize}

反面例子也说两个：

\begin{enumerate}

\item 用 \texttt{\#derive} 掩盖设计问题。 如果一个类型的相等语义需要"看起来相等就算相等"，派生逐字段比较会给出错误答案——这种情况该手写 \texttt{Equal}，而不是派生。

\item 用 \texttt{\#No\allowbreak{}Add\allowbreak{}Replace} 收拾本该拆开的类型。 如果 \texttt{a + 5} 和 \texttt{5 + a} 语义完全不同，往往说明这不该叫"加法"。先想想是不是该换成具名方法：\texttt{append(\allowbreak{}5)\allowbreak{}} 比 \texttt{operator+} 诚实。

\end{enumerate}

\section{边界与验证}

\begin{itemize}

\item 指令只修饰声明，不能出现在语句或表达式里；写在别处是语法错误。

\item \texttt{\#No<\allowbreak{}接口>\allowbreak{}Replace} 与 \texttt{\#No\allowbreak{}Right<\allowbreak{}接口>\allowbreak{}} 可以连写多条，顺序不影响结果。

\item 运算符开关只影响自动推导：直接写出的运算符、显式声明的 \texttt{Right\allowbreak{}Xxx} 都不受约束。

\item \texttt{Closeable} 自动等价于"不可复制"，无需再写 \texttt{\#Non\allowbreak{}Copyable}；\texttt{opcode} 的不可复制由声明体里有没有 \texttt{copy} 决定，\texttt{\#Non\allowbreak{}Copyable} 是可读性标记。

\item 可运行验证见 \texttt{example/\allowbreak{}samples/\allowbreak{}derive.\allowbreak{}chl}、\texttt{no\_replace.\allowbreak{}chl}、\texttt{closeable\_move.\allowbreak{}chl}、\texttt{operator\_right\_direction.\allowbreak{}chl}；报错例见 \texttt{example/\allowbreak{}samples-error/\allowbreak{}operator\_no\_add\_replace.\allowbreak{}chl}、\texttt{operator\_right\_forbidden.\allowbreak{}chl}、\texttt{closeable\_copy.\allowbreak{}chl}。

\end{itemize}

\section{小结}

{\small\setlength{\tabcolsep}{4pt}
\begin{longtable}{>{\raggedright\arraybackslash}p{\dimexpr 0.229\linewidth-2\tabcolsep\relax}>{\raggedright\arraybackslash}p{\dimexpr 0.751\linewidth-2\tabcolsep\relax}}
\toprule
你会了 & 关键点 \\
\midrule
\endhead
编译器指令是什么 & \texttt{\#} 开头、写在声明前、只给编译器看，不留到运行时 \\
\texttt{\#Non\allowbreak{}Copyable} & \texttt{opcode} 不绑定 \texttt{copy} \ensuremath{\rightarrow} 禁止复制，只能 \texttt{move(\allowbreak{}\&x)\allowbreak{}} \\
\texttt{Closeable} 的连带效果 & 实现即自动不可复制，\texttt{close(\allowbreak{})\allowbreak{}} 不会跑两遍 \\
\texttt{\#derive(\allowbreak{}.\allowbreak{}.\allowbreak{}.\allowbreak{})\allowbreak{}} & 白名单 Equal / Hash / Clone，嵌套字段要有协议，不能派生 Drop / Closeable \\
\texttt{\#No<\allowbreak{}接口>\allowbreak{}Replace} & 关掉该接口的全部自动推导（换边 / 对偶 / 取反） \\
\texttt{\#No\allowbreak{}Right<\allowbreak{}接口>\allowbreak{}} & 只关右向，左向与显式 \texttt{Right\allowbreak{}Xxx} 照旧 \\
报错怎么读 & 文案直接点出是哪条指令挡的，并给出两条出路 \\
泛型 & 指令随模板带走，实例化时贴回每个实例 \\
\bottomrule
\end{longtable}
}

\section{下一步}

到这儿，语言里所有"编译期能替你做的事"都讲完了：宏、注解、编译器指令。你已经知道什么时候该让编译器代劳，什么时候该自己动手。

最后一章是一份完整的实战项目：把结构体、反射、JSON 和 SQL 拼成一个能存数据的小系统。这一章不再介绍新语法，而是把前面二十几章的东西用一遍——你会发现真正花时间的地方，往往不是语法，而是"哪些边界该写进类型、哪些开关该关掉"。

翻到\hyperref[ch:24]{第 24 章 学生信息管理系统}。

\section{习题}

前两题动手验证指令的边界，第三题读报错，最后一题为综合挑战。

\begin{enumerate}

\item 给一个只实现 \texttt{Add<\allowbreak{}int64,\allowbreak{}int64>\allowbreak{}} 的结构体加上 \texttt{\#No\allowbreak{}Add\allowbreak{}Replace}，写 \texttt{5 + a} 并观察报错；把指令换成 \texttt{\#No\allowbreak{}Right\allowbreak{}Add}，再写一次，比较两条文案的差别。

\item 写一个 \texttt{\#derive(\allowbreak{}Equal,\allowbreak{} Hash,\allowbreak{} Clone)\allowbreak{}} 的结构体，字段里放一个自己实现 \texttt{Equal} 的嵌套结构体，确认派生能通过检查；把嵌套类型换成一个裸数组，观察报错。

\item 解释为什么 \texttt{\#derive(\allowbreak{}Drop)\allowbreak{}} 会被拒绝，而不是"生成一个空 drop"。

\item 挑战：设计一个"链式文本"类型，使 \texttt{chain + "x"} 合法、\texttt{"x" + chain} 编译报错，并说明为什么这里不需要 \texttt{Right\allowbreak{}Add}。

\end{enumerate}

\subsection*{参考解答与验收要点}

\begin{enumerate}

\item 前者文案是"运算符 + 的自动推导被 \#NoAddReplace 关闭"，后者是"右向被编译器指令关闭"；两者都保留了左向可用的行为。

\item 嵌套结构体需要自己实现或派生 \texttt{Equal} / \texttt{Hash} / \texttt{Clone}；裸数组没有约定语义，派生阶段会明确拒绝，而不是按字节比较。

\item \texttt{Drop} 是析构协议，值语义已经对字段做递归清理；再生成一个 \texttt{drop(\allowbreak{})\allowbreak{}} 会与自动清理叠加，等于析构两次。派生只处理"可以按字段机械推出"的协议。

\item 用 \texttt{\#No\allowbreak{}Right\allowbreak{}Add} + 左接口 \texttt{Add}；不需要 \texttt{Right\allowbreak{}Add}，因为右向语义根本不该存在。

\end{enumerate}

\section{番外：一张便条的权力}

改编自《聊斋志异·王者》。

某县衙门的库房里有一本账簿，凡入库的物件都要先登记。登记簿上原本没有"不得借出"这一栏，差役便自作主张：凡是不好搬的东西，就把借条贴在箱子上，权当它是规矩。

新来的主簿翻开库房，一眼看出问题：借条是写给取东西的人看的，规矩是写给管库房的人 看的。他把借条一张张揭下来，换成簿子上正经的条目——哪几类物件不得出库、哪几类出库要两个人签字、哪几类本来就该就地封存。从此取东西的人不必再猜哪张条子算数。

ChaLi 的编译器指令就是这簿子上的条目：它写在声明上，交给编译器读，读的人都一样、位置都在同一处。注解才是那张贴给外人看的借条，它不改变库房的规矩，只是告诉来办事的人里面装的是什么。

\subsection*{技术对照}

指令与注解的关键差别在读取者：指令由词法与语义层读取，参与类型检查；注解由反射在运行期读取，不改变语义。把"禁止某条推导"这类声明级约束写成标记接口，等于把规矩写进借条里。
